\scanpage{151}
\mspage{100}

L'hyp. sur $V',V'',V'\cap V''$ implique que les trois sommets du carré autres que $\underline{\unsure{\pi|\Pi}}_0(V/w)$ sont $\unsure{W|\mathcal W}$-asph, donc par $\mathrm W(7\mathrm{bis})$ il en est de $=$ de ce dernier, qed.

% slanted
\margin{Plus généralement, \deleted{\unsure{si}} si $I$ est une catégorie \uline{ordonnée} \uline{finie}, ayant un objet initial, alors}

\uline{Corollaire 2} Moyennant $\mathrm W(7\mathrm{bis})$, la catégorie

\diagram{\begin{tikzcd}[ampersand replacement=\&,column sep=small,row sep=small]
\& \& a\\
\unsure{\Psi|\Phi}={} \& 0 \arrow[ur] \arrow[dr] \&\\
\& \& b
\end{tikzcd}}

\diagram{\begin{tikzcd}[ampersand replacement=\&,column sep=small]
\text{(ou }0 \& 1\arrow[l]\arrow[r] \& 2\text{)}
\end{tikzcd}}

est tot. $W$-asphérique.

\inserted{\uuline{Faux}, \deleted{\unsure{$\Psi$}} n'est \uline{pas} \deleted{tot.} $W$-asphérique}

Qu'on applique le cor. 1 à $\Psi$ et à ses deux sous-catégories ouvertes, $\Psi_{/a}$, $\Psi_{/b}$ (isomorphes à $\Delta^1$), et à leur intersection $\{0\}\simeq\Delta^0=e$. Comme ces trois catégories ont un objet final (contrairement à $\Psi$ elle-même, qui n'a qu'un objet initial \unsure{0}), elles \deleted{sont} tot. $W$-asphériques. Le corollaire s'applique, l'existence des rétractions \inserted{$r'_0,r''_0$} étant ici évidente.

% slanted
\margin{\unsure{faux} : il y a les cat. filtrantes (si $W$ satisfait \unsure{à} l'axiome des \unsure{limites}). \deleted{Plus généralement les catégories telles \unsure{que} $X^{\mathrm o}$ ou bien $X$ \unsure{soit|tot.} \unsure{asph.}} \deleted{\unsure{i.e.} $X\to X\times X$ soit \unsure{asphérique} \unsure{(cf.} \unsure{chap.} \unsure{Théo)}, p. ex.} \deleted{$X=\Delta^0$, \unsure{ou} \unsure{qui} ont pour \unsure{but} \unsure{finissant},} \deleted{cf \unsure{X} \unsure{19} (p. 35).} Je conjecture \unsure{précisément|pourtant} que ce sont les seules !}

C'est surtout la partie b) \inserted{(rectifiée)} de la prop 9 \unsure{qui donne|démontre} l'existence de catégories tot. $W$-asphériques autres que celles ayant un objet final. Il est vrai qu'en sens inverse, il faudrait quand même s'assurer qu'un objet $X$ $W$-asph. de $\mathrm{Cat}$ n'est pas \uline{toujours} tot. $W$-asphériques ! Le \unsure{contrexemple|contre-exemple} n'est pas tellement évident finalement --
% archive sheet number: 151
